	\documentclass[14pt]{extreport}
\usepackage{extsizes}
	\usepackage[frenchb]{babel}
	\usepackage[utf8]{inputenc}  
	\usepackage[T1]{fontenc}
	\usepackage{amssymb}
	\usepackage[mathscr]{euscript}
	\usepackage{stmaryrd}
	\usepackage{amsmath}
	\usepackage{tikz}
	\usepackage[all,cmtip]{xy}
	\usepackage{amsthm}
	\usepackage{varioref}
	\usepackage[ margin=1in]{geometry}
	\geometry{a4paper}
	\usepackage{lmodern}
	\usepackage{hyperref}
	\usepackage{array}
	\usepackage{float}
	\usepackage{easytable}
	 \usepackage{fancyhdr}\usepackage{longtable}
	 \usetikzlibrary{shapes.misc}
	 \newcommand\ang[1]{$#1{}^o$}
\newlength{\taillecellule}
\setlength{\taillecellule}{2cm}
\newcolumntype{C}{@{}>{\centering\arraybackslash}p{\taillecellule}@{}}

\renewcommand{\theenumi}{\alph{enumi})}
\usepackage{pstricks,multido}
\usepackage{arrayjob}
\usepackage{calc,xlop}
\tikzset{cross/.style={cross out, draw=black, minimum size=2*(#1-\pgflinewidth), inner sep=0pt, outer sep=0pt},
%default radius will be 1pt. 
cross/.default={1pt}}

	\newcounter{n}
	\setcounter{n}{1} 
	\newcommand{\cm}{\textbf{ cm}}
	\newcommand\ds\displaystyle
	\renewcommand{\it}{\item[$ \mathbf{\bullet\ \then}.$]\stepcounter{n} }
	\newenvironment{calcul}{\item[$\mathbf{\then}.$] \stepcounter{n}\hfil $\displaystyle  }{.$\hfil   }
	\pagestyle{fancy}
	\theoremstyle{plain}
	\fancyfoot[C]{\empty} 
	\fancyhead[L]{Bilan}
	\fancyhead[R]{Année de 3e}
	
	
	
	\title{Bilan année de 3e}
	\date{}
	\begin{document}

\subsection*{Avant-propos}
Je vous ai rédigé ces exercices (\url{https://leturcqd.github.io/ressources/3e/R\%C3\%A9visions3e/R\%C3\%A9visions3e.pdf}) comme une base de travail pour ceux qui veulent se renforcer un peu avant le lycée, et aussi pour ceux qui sont en recherche de problèmes plus intéressants. \\La première partie comporte des exercices de calcul pur sur le programme de collège : quatre opérations sur les relatifs et les fractions, manipulations des signes et des parenthèses, calcul littéral (réduction, distributivité)\footnote{Cela manque de factorisations, il faudra que je vous en rajoute en v2.}, résolution d'équations.
\\ La seconde partie comporte un mélange de problèmes plus autonomes utilisant des choses que vous avez vues à un moment ou à un autre : problèmes de fractions, de vitesses, de géométrie des polygones, énigmes numériques, puis des problèmes plus uniques et originaux à partir de la question w). \\
Si j'en trouve le temps, je vous ajouterai des corrections de tout cela dans quelques temps au lien suivant : \url{https://leturcqd.github.io/ressources/3e/R\%C3\%A9visions3e/R\%C3\%A9visions3ec.pdf}
\subsection*{Révision de calcul}  \noindent
Effectuer les additions suivantes : 
\begin{itemize}
\it $\left(-\frac14\right)+\left(-\frac32\right)+\left(+\frac52\right)$.
\it $\left(-\frac56\right)+\left(+\frac49\right)+\left(-\frac73\right)$.
\it $(-0,75) + (+4) + \left(-\frac45\right)$.
\it $\left(+\frac47\right)+\left(-\frac{13}{21}\right)+\left(+\frac56\right)$.
\it $(-12)+(-11,57)+(+9,87)$. 
\it $ (-4,51)+\left(-\frac45\right)+\left(+\frac7{20}\right)$.
\end{itemize}
Calculer la valeur des sommes algébriques suivantes : 
\begin{itemize}
\it $(-10)+(-15)-(-30)+(+7)-(+12)$.
\it $(-7)-(-9)-(+17)+(-41)+(+1)$.
\it $(+49)-(-63)+(-3)-(+4)$.
\end{itemize}
Calculer la valeur des sommes algébriques suivantes : 
\begin{itemize}
\it $(-0,7)-(-0,9)+(-13,5)-(+4)$
\it $\left(-\frac12\right)+\left(-\frac34\right) - \left(+\frac25\right)-(-1)$.
\it $\left(-\frac5{11}\right)+\left(-\frac7{22}\right)
-\left(-\frac4{33}\right) + (-1) -(+0,5)$.
\end{itemize}
Calculer la valeur des sommes algébriques suivantes : 
\begin{itemize}
\it $7-10+14-13-21+3$.
\it $0,75+4,7-0,7-0,5-1$.
\it $\frac23-\frac34+\frac5{12}-\frac56 - 3$. 
\end{itemize}
Effectuer les opérations suivantes : 
\begin{itemize}
\it $[(5-9)+(3-5)] - [(7+3-5) - (7-10)]$.
\it $[12-(14-5+0,75)] + [-15 + (3,25-2)]$.
\it $\left[1 - \left(\frac25-\frac43\right)\right] - 
\left[ \left(\frac45-1\right) - \left(\frac75+2\right)\right]$. 
\end{itemize}
Supprimer les parenthèses et les crochets dans les expressions suivantes : 
\begin{itemize}
\it $(a-b+c) - (d- e - f) + (b-a)$. 
\it $[(a-b)-(a-5)] + [b-7-(a-3)]$. 
\it $[12-(a-b)+6] - [15+ (b-a-13)]$. 
\end{itemize}

Effectuer les produits suivants : 
\begin{itemize}
\it $(-3)(+7)(-5)(-4)$.
\it $(+1,7)(-2,5)(-3)(-1,9)$.
\it $(-1)(+1)(-1)(+1)(-1)(-1)$.
\it $\left(-\frac35\right)(+3)\left(-\frac53\right)$.
\it $\left(+\frac12\right)\left(-\frac14\right)\left(-\frac48\right)$.
\it $\left(+\frac12\right)\left(-\frac6{11}\right)\left(-\frac78\right)\left(-\frac{11}{15}\right)$.
\end{itemize}
Effectuer les divisions suivantes : 
\begin{itemize}
\it $\displaystyle(+15):\left(-\frac23\right)$.
\it $\displaystyle(-7,5):\left(-\frac45\right)$.
\it $\displaystyle(-1):\left(+\frac34\right)$.
\it $\displaystyle\left(-\frac45\right):(+1,4)$.
\it $\displaystyle\left(+\frac34\right):\left(-\frac{15}8\right)$.
\it $\displaystyle\left(-\frac73\right):\left(+\frac53\right)$.
\it $\displaystyle(+1):\left(+\frac23\right)$.
\it $\displaystyle\left(-\frac{8}{25}\right):\left(-\frac{14}5\right)$.
\it $\displaystyle\left(+\frac13\right):\left(-\frac14\right)$.
\end{itemize}
Effectuer les opérations suivantes : 
\begin{itemize}
\it $\displaystyle\frac{-3+5-6}{2+7-1}\times\frac{4-3+1}{2-3-1}\times \frac{-5-7}{10}$.
\it $\displaystyle\frac{\frac12-\frac23}{\frac34-1}\times \frac{\frac56+\frac13}{1-\frac45}\times \frac{-\frac78+1}{\frac32}.$
\it Calculer : \[ (-11)^3; \phantom{miou}(+9)^2; \phantom{miou}(-10)^5; \phantom{miou}(+7)^4; \phantom{miou}(-13)^4.\]
\it Calculer : \[ \left(-\frac12\right)^5; \phantom{miou}\left(+\frac23\right)^2; \phantom{miou}\left(-\frac34\right)^3; \phantom{miou}(+1,4)^2; \phantom{miou}(-0,25)^2.\]
\it Calculer : \[ 5^3\times 5^2; \phantom{miou}(-7)\times(-7)^2; \phantom{miou}\left(\frac23\right)\times \left(\frac23\right)^3; \phantom{miou}
\left[\left(-\frac14\right)^2\right]^2\left[(-2)^2\right]^3\]
\end{itemize}
Réduire les expressions suivantes: 
\begin{itemize}
\it $-\frac32 x + \frac54 x-3x^2+\frac{x}6-\frac52 x^2+5+4x^2$.
\it $\frac32 x^2+xy+y^2-2xy+\frac{x^2}3-\frac32x^2$.
\it $4a^2-\frac23a-\frac35a^2+\frac13a-5a-\frac2{15}a^2$.
\it $3x^2+\frac45 - \frac53x - 2x^2 -\frac35x^3 + 4 - 2 x^2 + 7x$. 
\it $4x^2 - \frac72 + \frac35x - \frac52x^2 + \frac43 x^3 - 5 + \frac32 x^3 + 7 - 2x$. 
\it $\frac25 a^2b + 3a^3 -4ab^2 + \frac52 a^2b + \frac72b^3 - b^3 + 2ab^2$. 
\end{itemize}
Réduire : \begin{itemize}
\it $(3x-5)+[2x-5 - (3x-2y+4) - (4x-3y-9)]$.
\it $(2x-5y+7) - [(3x+2y-3)-(4x+4y-2)]-[2x-(3y+4)]$. 
\it $[(x-2y+5) - (3x+2y+7)]-[(2x+3)-(4y-2)]$. 
\end{itemize}
Effectuer les produits suivants : 
\begin{itemize}
\it $(3a^2b^3)\left(\frac23ab^5\right)$.
\it $\left(\frac45a^3b^2c\right)\left(-\frac34abc^4\right)$. 
\it $\left(\frac47a^2xy^3\right)\left(-\frac52a^3y^4\right)$.
\it $\left(-\frac34x^2y\right)\left(+\frac35a^3y^5\right)$.
\it $\left(\frac94a^4x^2y^3\right)\left(-\frac43ax^2\right)$. 
\it $\left(\frac{14}3a^2b^3x\right)\left(-\frac67a^2b^5\right)$. 
\it $\left(-\frac72ax^2y\right)\left(-\frac8{15}b^3xy^2\right)\left(\frac5{21}abx^3\right)$.
\it $\left(-\frac23xy^2\right)^2(-4x^2y)$.
\it $\left(\frac5{12}a^4b^2x\right)\left(-\frac27ax^2y^3\right)\left(-\frac{14}5b^2xy^4\right)$.
\it $\left(\frac35x^2y\right)^3\left(-\frac54xy\right)$.
\end{itemize}
Effectuer les produits suivants.
\begin{itemize}
\it $(2x-7)(-3x+2)$.
\it $(4x^5+7-2x^3)(x^3-2x)$.
\it $(5x^3-2x)(3x-4x^2)$
\it $(2x-7x^2+5x^3)(3x-5x^2+8)$.
\it $\left(-2x+\frac32\right)(4x+3)$.
\it $\left(\frac83x-\frac32x^2+5\right)(4x^3-5x^2+7)$.
\it $(7x^4-2x^3+4x^2)(3x^2-5)$.
\it $(2x^2-4x^3)(x^3-2x)$.
\it $(2x^2-4+2x)(x^2+5-2x)$.
\it $\left(\frac54x^3-2x+\frac12\right)\left(\frac72x^3-\frac23x+x^2\right)$.
\end{itemize}
 Calculer les produits suivants : 
\begin{itemize}
\it $(2x+3)(3x+2)(x-4)$.
\it $(5x-1)(2x+3)(7+4x)$.
\it $(3x^2-1)(x+1)(x-1)$. \end{itemize}
Développer et réduire : 
\begin{itemize}
\it $5(3a^2-4b^3)-[9(2a^2-b^3)-2(a^2-5b^3)]$.
\it $3a^2(2b-1)-[2a^2(5b-3)-2b(3a^2+1)]$.
\it $(2a+5b)(3a-2b)-(2a-1)(3a+2b)-(a-2b)(5b-1)$.
\it $(2x-3y)(5x-2y)-(3x-2y)(2x+1)-(5x-y)(3y+1)$.
\it $(ax^2-b)(ax^2-2b)+3b(ax^2-b)+b(b-1)$.\end{itemize}
Résoudre les équations suivantes : 
\begin{itemize}
\it $5(2x-3) - 4(5x-7)=19-2(x+11)$. 
\it $4(x+3)-7x+17 = 8(5x-1)+166$. 
\it $17-14(x+1)= 13-4(x+1)-5(x-3)$. 
\it $5x+3,5+(3x-4) = 7x-3(x-0,5)$. 
\it $7(4x+3)-4(x-1)=15(x+0,75)+7$. 
\it $17x+15(x-1)=-1-14(3x+1)$. 
\end{itemize}
Résoudre les équations suivantes (après développement, les termes en $x^2$ ou $x^3$ se simplifient) : 
\begin{itemize}
\it $(x-1)^2+(x+3)^2= 2(x-2)(x+1)+38$.
\it $5(x^2-2x-1)+2(3x-2)=5(x+1)^2$.
\it $(9x+1)(x-2)=(3x+4)(3x-5)$.\end{itemize}
Résoudre les équations suivantes (on se ramènera au cas entier en multipliant par un nombre opportun).
\begin{itemize}
\it $\frac{x+5}4 - \frac{x-3}6 = \frac{x}3$. 
\it $\frac{3x-7}2 + \frac{x+1}3 = -16$. 
\it $x-\frac{x+1}3 = \frac{2x+1}5$. 
\it $\frac{7-3x}{12} + \frac34 = 2(x-2) + \frac{5(5-2x)}6$. 
\it $\frac{x}5 - \frac{3x-1}6 + \frac{3-x}4 = 0$.
\it $\frac{3(x+3)}4 + \frac12 = \frac{5x+9}3 - \frac{7x-9}4$. 
\it $\frac{2x-7}5 + \frac{x+11}2 = -4$. 
\it $\frac{2x-3}3 - \frac{x-3}6 = \frac{4x+3}4 - 17$. 
\it $\frac{5x-3}4 - \frac{7x-5}9 = \frac{x+19}6$. 
\it $\frac{5x+1}8-\frac{x-1}3 = \frac{4(2x-3)}9$. 
\it $\frac{2x-1}3 - \frac{5x+2}7 = x +13$. 
\it $\frac{8x+2}5 - \frac{x-11}7 = \frac{5x-3}2 - \frac{3x-1}4$.\end{itemize}
Résoudre les équations :
\begin{itemize}
\it $(x-1)(x+2)(x-3)=0$. 
\it $(x-3)(x-4)(x-5)=0$. 
\it $(2x+1)(x+1)(4x-3) = 0$.
\it $(2x+1)(x+4)(3x+1)=0$. 
\it $x(5x+1)(4x-3)(3x-4)=0$. 
\it $5x(3x-7)=0$.   
\it $x^2-3x=0$. 
\it $5x^2+8x=0$. 
\it $4x^2-\frac{7x}3=0$. 
\end{itemize}
	
    \subsection*{Problèmes de réflexion plus avancés}
\begin{enumerate}
\item Quel nombre faut-il ajouter aux deux termes de la fraction $\frac7{13}$ pour qu'elle devienne égale à $\frac23$ ? 
 \item Quel nombre faut-il retrancher aux deux termes de la fraction $\frac{17}{23}$ pour obtenir une fraction égale à $\frac58$ ? 
 \item Quel nombre faut-il ajouter aux deux termes de la fraction $\frac58$ et retrancher aux deux termes de la fraction $\frac34$ pour obtenir deux fractions égales ? 
 \item Une personne dispose de deux heures pour effectuer une promenade. Elle part en tramway à la vitesse moyenne de $12$ km à l'heure et revient à pied, à la vitesse moyenne de $4$ km à l'heure. À quelle distance du point de départ devra-t-elle quitter le tramway ? 
 \item Un épicier achète un certain nombre de bouteilles de vin pour $225$ euros. En vendant la bouteille $1,65$ euros, il gagnerait autant que ce qu'il perdrait s'il la vendait $1,35$ euros. Quel est le nombre des bouteilles de vins ? 
 
 \item On considère un polygone régulier à $6$ côtés.
 \begin{enumerate}
 \item Réaliser une figure.
 \item Relier tous les sommets au centre $O$. 
 \item Montrer que tous les triangles ainsi formés ont leurs angles tous égaux à $60^o$. 
 \item En déduire que les angles d'un hexagone régulier mesurent $120$ degrés.
\end{enumerate}   

 \item On considère un polygone régulier à $5$ côtés.
 \begin{enumerate}
 \item Réaliser une figure.
 \item Relier tous les sommets au centre $O$. 
 \item Montrer que tous les triangles isocèles ainsi formés ont leurs angles à la base tous égaux à $54^o$. 
 \item En déduire que les angles d'un pentagone régulier mesurent $108^o$ degrés.
\end{enumerate}   
\item Utiliser les deux exercices précédents pour généraliser à un polygone régulier avec un nombre quelconque $N$ de côtés. On montrera ainsi que les angles d'une telle figure valent tous $\frac{N-2}N\times 180^o$.
\item On trace deux points $A$ et $B$ et une droite $(d)$ sur une feuille. Expliquer comment construire le cercle passant par $A$ et $B$ dont le centre appartient à $(d)$. Existe-t-il toujours ? Y en a-t-il toujours un seul ?
\item On place trois points $A$, $B$, $C$ sur une feuille. À quelle condition passe-t-il un cercle par ces trois points ? Comment le construire ? 
\item On considère un quadrilatère convexe $ABCD$ dans lequel la médiatrice de 
$[AB]$ est également la médiatrice de $[CD]$. Soit $O$ le point où cette médiatrice coupe la médiatrice de $[BC]$. \begin{enumerate}
\item Montrer que les quatre points $A$, $B$, $C$ et $D$ sont situés sur un même cercle de centre $O$. 
\item Démontrer que les médiatrices de $[AD]$, $[AC]$ et $[BD]$ passent également par $O$. 
\end{enumerate}
\item Tracer un parallélogramme $ABCD$, et les bissectrices de ses quatre angles. Elles se croisent intérieurement au parallélogramme en quatre points $E$, $F$, $G$ et $H$. Que dire du quadrilatère $EFGH$ formé ? Le démontrer. 
\item On trace un triangle $ABC$ et on place les milieux $D$, $E$, $F$ de ses trois côtés. Montrer que le triangle $DEF$ est une réduction du triangle $ABC$ dont on déterminera le coefficient. 
\item Dans le même contexte que l'exercice précédent, on relie chaque milieu au sommet opposé, et on note $G$ le point d'intersection des trois segments ainsi formés. Montrer que l'homothétie de centre $G$ et de rapport $-2$ envoie chaque sommet du triangle $DEF$ sur un sommet du triangle $ABC$. 
\item Dans un triangle $ABC$, on trace la hauteur issue de $A$ qui rencontre\footnote{Attention, la hauteur coupe le segment, c'est une vraie hypothèse : si elle rencontrait $(BC)$ hors du segment, le résultat serait faux !} le segment $[BC]$ en $H$. On suppose $AB^2 = BH \times BC$. Montrer que $ABC$ est rectangle en $A$. 
\item On considère un quadrilatère $ABCD$. Montrer que les milieux des côtés forment un  parallélogramme.
 \item Trouver trois nombres entiers consécutifs (consécutifs : qui se suivent) dont la somme est égale à $57$. 
 \item La somme de cinq nombres impairs consécutifs est 85. Quels sont ces nombres ? 
 \item Un père a 29 ans ; son fils a 5 ans ; dans combien d'années l'âge du père sera-t-il le triple de l'âge du fils ? 
 \item Un père a 41 ans ; ses trois enfants sont âgés de 7, 9 et 13 ans. Dans combien d'années l'âge du père sera-t-il égal à la somme des âges de ses enfants ? 
 \item Trouver un nombre dont la somme des quotients par 5, 7 et 9 soit égale à 429. 
 \item Développer les produits suivants : 
 \[ (x-1)(x+1)\]
 \[(x-1)(x^2+x+1)\]
 \[(x-1)(x^3+x^2+x+1)\]
 En déduire une formule générale. En déduire que 
 \[ x^n+ x^{n-1}+\ldots + x^2 + x + 1 = \frac{x^{n+1}-1}{x-1}\]
 \item \  [Utiliser l'exercice précédent] Conte classique : un sage demande à un roi de le payer en mettant une pièce sur la première case d'un échiquier, puis deux pièces sur la seconde, puis quatre pièces, et ainsi de suite sur les soixante-quatre cases en doublant à chaque nouvelle case le montant. Déterminer la somme reçue. En utilisant l'approximation $2^{10} \sim 1000$, en donner un ordre de grandeur. 
 \item Un cercle délimite deux régions du plan : l'intérieur et l'extérieur. Deux cercles peuvent en déterminer quatre. S'en convaincre par un dessin. Que dire du nombre maximal de régions délimitées par trois cercles ? Que se passe-t-il pour quatre ? Trouver une formule générale.
 \item Si on coupe un disque par une corde, on délimite deux régions. Une deuxième corde délimite alors quatre régions. Généraliser. 
 \item Un nombre entier $N$ est dit \emph{pratique} si chaque entier entre $1$ et $N$ peut être écrit comme une somme de diviseurs de $N$ (sans répétition). Par exemple, $6$ est pratique car ses diviseurs sont $1$, $2$, $3$ et $6$ et que : 
 \[ 1 = 1\] 
 \[ 2 = 2 \]
 \[ 3 = 1 + 2 \]
 \[ 4 = 1+3 \text{ (mais 2+2 ne serait pas autorisé)}\]
 \[ 5 = 2 + 3 \]
 \[ 6 = 6\]
 Déterminer les nombres pratiques jusqu'à $30$. \\
 Montrer que toute puissance de $2$ est pratique (Penser à l'écriture en base $2$). 
\end{enumerate} 
 
\end{document}